 この章では、本システムの機能やそのアルゴリズムについて説明する。

クラス
html
AnimationLine
AnimationEllipse
AnimationArc
AnimationText
AnimationPolygon
SlideAnimation
figure
FigureAnimation
AnimationSetting
position

\section{アニメーション描画領域}
 fsxファイルにおいてソースコード\ref{animation_fsx}のような記述を行うことで、fsファイルに定義されたソースコード\ref{animation_fs}が呼び出され、アニメーション描画領域が生成される。
  \begin{lstlisting}[caption={アニメーション描画領域を生成するfsxファイル},label={animation_fsx},breaklines=true,breakatwhitespace=false]
html.animation outputdir "PresentationSlide_nishimura" {sX=800; sY=580; mX=1040; mY=450; backgroundColor="#bbeeff"} p (960,450) <| fun (f,p) ->
\end{lstlisting}
  \begin{lstlisting}[caption={アニメーション描画領域の生成処理を定義したfsファイル},label={animation_fs},breaklines=true,breakatwhitespace=false]
static member animation (outputdir:string) (projectname:string) (s:ViewBoxStyle) (p:position) (buttonX:int,buttonY:int) code =
    figcounter <- figcounter + 1
    let f = FigureAnimation(wrBody,wrJS,figcounter,outputdir,projectname,s.mX,s.mY,s.sX,s.sY)
    wrJS.Write "function repeatSeq(fn, interval, Nt, onComplete)"
    //時間制御をおこなう処理　（省略）
    write ("<svg viewBox=\"0 0 "+s.sX.ToString()+" "+s.sY.ToString()+"\" ")
    //SVG要素を生成および属性設定処理　（省略）
    write "</svg>"
    let asc = nextAnimationGroup()
    let fnameStart = f.jsStartControll asc
    let fnameReset = f.jsResetControll asc
    addAnimationButton (fnameStart,fnameReset,buttonX,buttonY)
\end{lstlisting}
ソースコード\ref{animation_fs}では、htmlクラスに静的メンバ関数として描画領域生成処理を定義している。本関数は五つの引数を取る。
第一引数は、出力ファイルの保存先を指定する。第二引数は出力するファイル名を指定する。
第三引数はアニメーション描画領域の設定を行う。第四引数はアニメーション描画領域の描画位置を指定する。第五引数はアニメーションボタンの位置を指定する。
\par
4～5行目では、アニメーションの時間制御を行うJavaScriptファイルを自動生成する処理を実装している。
6～8行目では、アニメーションの描画領域を設定するHTMLファイルを自動生成する処理を実装している。これらの処理の詳細は、\ref{subsec:javascript}節および\ref{subsec:html}節において述べる。

\subsection{JavaScriptによる時間制御処理}\label{subsec:javascript}
 ソースコード\ref{repeatSeq}は、アニメーション生成における時間制御処理を行う処理を実装したコードである。
 \begin{lstlisting}[caption={時間制御処理},label={repeatSeq},breaklines=true,breakatwhitespace=false]
wrJS.Write "function repeatSeq(fn, interval, Nt, onComplete)"
wrJS.Write "{"
wrJS.Write "    let t = 0;"
wrJS.Write "    function run()"
wrJS.Write "    {"
wrJS.Write "        if (t < Nt)"
wrJS.Write "        {"
wrJS.Write "            fn(t);"
wrJS.Write "            t++;"
wrJS.Write "            setTimeout(run, interval);"
wrJS.Write "        }"
wrJS.Write "        else"
wrJS.Write "        {"
wrJS.Write "            onComplete();"
wrJS.Write "        }"
wrJS.Write "    }"
wrJS.Write "    run();"
wrJS.Write "}"
\end{lstlisting}
まず、1行目にて
\subsection{HTMLによる描画領域設定}\label{subsec:html}
\subsection{ボタン表示位置の設定}
\section{ボタン効果}
\subsection{startボタン}
\subsection{resetボタン}
\section{図形}
\subsection{生成と表示}
\subsection{スタイル制御}
\subsection{時間に応じた変形・移動}
\subsection{複数図形の同時制御}
\section{画像}
\section{テキスト}
\section{アニメーション図説フレームワークの機能}